% Copyright (c) 2026 Eugene Joseph M. Ragasa.
% SPDX-License-Identifier: Apache-2.0
% Author: Eugene Joseph M. Ragasa
% Formal proof completed: 2026-10-20
\documentclass[11pt]{article}
\usepackage[T1]{fontenc}
\usepackage{amsmath,amssymb,booktabs,tabularx}
\usepackage{listings}
\usepackage[margin=1in]{geometry}
\usepackage{xurl}
\usepackage[hidelinks]{hyperref}
\setlength{\emergencystretch}{3em}
\lstset{basicstyle=\ttfamily\small,breaklines=true,columns=fullflexible}
\title{Proof Dossier: Exact Lower-Band Projector Sandwich Identity}
\author{Eugene Joseph M. Ragasa}
\date{Formal proof completed: 2026-10-20}
\begin{document}\sloppy\maketitle
\section{Identity and status}
\begin{tabularx}{\textwidth}{@{}lX@{}}\toprule
Statement ID & \path{BG-PROJECTION-THM-006} \\
Status & proved; no admissions \\
Lean declaration & \path{ProjectKoios.BandGauge.Projection.twoBandLowerBandProjector_sandwich} \\
Owning module & \path{ProjectKoios/BandGauge/Projection/SpectralProjectionRange.lean} \\
Axioms & \texttt{propext}, \texttt{Classical.choice}, \texttt{Quot.sound} \\
\bottomrule\end{tabularx}

\section{Exact mathematical statement}
For every $k\in\mathbb R$ and real $\Delta,t$ satisfying $\Delta t\ne0$, the source
parent, dispersion, and lower-band projector satisfy
\[
 P_-(k)H(k)P_-(k)=\frac12\bigl(H(k)-d(k)I\bigr).
\]

\section{Source provenance and relationship}
This theorem \emph{formalizes} equation \texttt{eq:projected-range} in
\path{docs/publications/papers/Ragasa2026b/v2/manuscript.tex:328--334}, owner commit
\path{99a46cb59d0c6f27ada7d709f194996ef9619d9f}, SHA-256
\texttt{d8d5951d8d7ed75355b5d94f37fcbddc70c1466897a5770dbe9f877e60016daf}.

\section{Assumptions, conventions, and delta}
The exact source nondegeneracy $\Delta t\ne0$ is preserved. The projector, parent, and
positive dispersion are the declarations inventoried as
\path{BG-PROJECTION-DEF-006/009/010/011}. Equality is pointwise in real momentum and
exact over complex $2\times2$ matrices. There is no numerical approximation.

\section{Statement-specific prior art}
\begin{tabularx}{\textwidth}{@{}lXXX@{}}\toprule
Relation & Identity & Locator & Role \\
\midrule
formalizes & owner source equation & manuscript lines 328--334 & exact displayed identity \\
depends-on & \path{BG-PROJECTION-THM-001/005} & theorem manifest & parent square and projector specification \\
background-only & external identity unresolved & unverified & standard rank-two projector algebra; no novelty claim \\
\bottomrule\end{tabularx}

\section{Motivation and contribution boundary}
This identity connects the algebraic projector to the exact matrix function whose finite
Fourier support is tested next. It does not itself prove non-finiteness or a universal
projection theorem.

\section{Proof schematic and argument}
Use $P_-=(I-H/d)/2$ and $H^2=d^2I$. Direct expansion gives
\[
 P_-HP_-=\frac14(I-H/d)H(I-H/d)
 =\frac12(H-dI).
\]
The gap hypothesis makes $d$ nonzero, so each cancellation is valid.

\section{Lean realization}
A private generic matrix lemma expands the sandwich using matrix/scalar multiplication,
substitutes the parent-square theorem, and normalizes the remaining field expression.
The source declaration instantiates it at each angle. Toolchain: Lean 4.35.0, Mathlib
\texttt{56a005855ea17e0b0958f73ec801ae140b0a1679}; no additional dependency.

\section{Trusted base and axiom report}
\begin{lstlisting}
#print axioms ProjectKoios.BandGauge.Projection.twoBandLowerBandProjector_sandwich
-- [propext, Classical.choice, Quot.sound]
\end{lstlisting}

\section{AI, automation, and comprehension record}
The implementation and prose were AI-assisted. Lean kernel checking validates the exact
identity. Field and ring normalization produced checked proof terms. The maintainer
reviewed the expansion and source correspondence. No independent expert review occurred.

\section{Scientific and evidence boundaries}
No runtime artifact, numerical eigensystem, material interpretation, hopping truncation,
or decay estimate is involved.

\section{Unresolved questions}
A general reusable theorem for arbitrary finite-dimensional involutive algebras may be
worth extracting later. Canonical external prior-art identity remains unresolved.

\section{Review and revision history}
\begin{tabularx}{\textwidth}{@{}lllX@{}}\toprule
Date & Revision & Process & Outcome \\
\midrule
2026-10-20 & initial & AI-assisted implementation; Lean kernel check & exact source sandwich identity proved; no independent expert review \\
\bottomrule\end{tabularx}
\end{document}
